% 雨果·斯坦豪斯（Hugo Steinhaus）（综述）
% license CCBYSA3
% type Wiki

本文根据 CC-BY-SA 协议转载翻译自维基百科\href{https://en.wikipedia.org/wiki/Hugo_Steinhaus}{相关文章}。

\begin{figure}[ht]
\centering
\includegraphics[width=6cm]{./figures/2db84a4d97edc48c.png}
\caption{雨果·斯坦豪斯（1968年）} \label{fig_YGstgs_1}
\end{figure}
雨果·迪奥尼济·斯坦豪斯（Hugo Dyonizy Steinhaus，英语发音：/ˈhjuːɡoʊ ˈstaɪnˌhaʊs/，波兰语发音：[ˈxuɡɔ ˈʂtajnxaws]，1887年1月14日－1972年2月25日）是一位波兰数学家和教育家。斯坦豪斯于1911年在哥廷根大学师从大卫·希尔伯特（David Hilbert）获得博士学位，后来成为利沃夫扬·卡齐米日大学（现今乌克兰利维夫）的教授，并在那里帮助建立了后来被称为“利沃夫数学学派”（Lwów School of Mathematics）的学术中心。他被誉为“发现”数学家斯特凡·巴拿赫（Stefan Banach）的人，与巴拿赫共同在泛函分析领域作出了重要贡献，特别是提出了著名的巴拿赫–斯坦豪斯定理（Banach–Steinhaus theorem）。

第二次世界大战后，斯坦豪斯在弗罗茨瓦夫大学（Wrocław University）数学系的建立以及战后波兰数学界的重建中发挥了重要作用。斯坦豪斯一生发表了约170篇科学论文与著作，他在数学的多个分支领域留下了深远影响，包括泛函分析、几何学、数理逻辑以及三角学。值得注意的是，他被认为是博弈论（game theory）和概率论（probability theory）的早期奠基者之一，为后续学者发展更系统的理论框架奠定了基础。
\subsection{早年生活与求学经历}
斯坦豪斯于1887年1月14日出生于奥匈帝国的亚斯沃（Jasło）一个有犹太血统的家庭。\(^\text{[1]}\)他的父亲博古斯瓦夫（Bogusław）是一名地方实业家，拥有一家砖厂，同时也是一位商人；他的母亲名为埃韦丽娜（Ewelina），娘家姓利普希茨（Lipschitz）。斯坦豪斯的叔叔伊格纳奇·斯坦豪斯（Ignacy Steinhaus）\(^\text{[2]}\)是一名律师，也是“波兰议会俱乐部”（Koło Polskie）的成员，并担任加利西亚和洛多梅里亚王国（Kingdom of Galicia and Lodomeria）地区议会——加利西亚议会（Galician Diet）的代表。
\begin{figure}[ht]
\centering
\includegraphics[width=8cm]{./figures/844c35d7f6c20eff.png}
\caption{1907年，斯坦豪斯于哥廷根（Göttingen）拍摄的合影中站在最右侧。} \label{fig_YGstgs_2}
\end{figure}
1905年，斯坦豪斯在亚斯沃文理中学（gymnasium in Jasło）完成学业。他的家人原本希望他成为一名工程师，但他却被抽象数学深深吸引，开始自学当时著名数学家的著作。同年，他进入伦贝格大学（今乌克兰利维夫大学，University of Lemberg）学习哲学与数学。\(^\text{[2]}\) 1906年，他转入哥廷根大学（Göttingen University）。\(^\text{[1]}\)在那里，他师从大卫·希尔伯特（David Hilbert），并于1911年获得博士学位，论文题目为《Dirichlet 原理的新应用》（德文原题：Neue Anwendungen des Dirichlet'schen Prinzips）。\(^\text{[3]}\)

第一次世界大战爆发后，斯坦豪斯回到波兰，加入约瑟夫·毕苏斯基（Józef Piłsudski）领导的波兰军团（Polish Legion）服役，战后定居于克拉科夫（Kraków）。\(^\text{[4]}\)斯坦豪斯是无神论者。\(^\text{[5]}\)
\subsection{学术生涯}
\subsubsection{两次世界大战之间的波兰时期}
在1916年至1917年期间，即波兰于1918年重新获得独立之前，斯坦豪斯曾在克拉科夫（Kraków）为内务部工作，服务于短暂存在的傀儡政权——波兰王国（Kingdom of Poland）。\(^\text{[6]}\)

1917年，他开始在伦贝格大学（University of Lemberg，后来的波兰扬·卡齐米日大学 Jan Kazimierz University）任教，并于1920年取得讲师资格（habilitation）。\(^\text{[6]}\) 1921年，他被任命为“副教授”（profesor nadzwyczajny），1925年晋升为“正教授”（profesor zwyczajny）。\(^\text{[6]}\)在此期间，他开设了当时仍属前沿领域的勒贝格积分理论（Lebesgue integration）课程——这是法国以外最早开设此类课程的大学之一。\(^\text{[3]}\)

在利沃夫（Lwów）任教期间，斯坦豪斯共同创立了著名的“利沃夫数学学派”（Lwów School of Mathematics）。\(^\text{[7]}\)他也是苏格兰咖啡馆（Scottish Café）中活跃的数学家圈成员之一，不过据斯坦尼斯瓦夫·乌拉姆（Stanisław Ulam）回忆，斯坦豪斯本人在这些聚会上通常更喜欢到街上那家更为高雅的茶馆小坐。\(^\text{[4]}\)
\subsubsection{第二次世界大战时期}
\begin{figure}[ht]
\centering
\includegraphics[width=6cm]{./figures/6c566e9ff83fe38a.png}
\caption{} \label{fig_YGstgs_3}
\end{figure}
1939年9月，纳粹德国与苏联分别入侵并占领波兰，以履行先前签署的《莫洛托夫–里宾特洛甫条约》。利沃夫（Lwów）最初落入苏联控制之下。斯坦豪斯曾考虑逃往匈牙利，但最终决定留在利沃夫。苏联方面对大学进行了重组，使其更具乌克兰特色，但他们任命了斯坦豪斯的学生斯特凡·巴拿赫（Stefan Banach）为数学系主任，斯坦豪斯也得以恢复教学工作。数学系师资队伍此时还因多位来自德占波兰的难民学者而得到加强。斯坦豪斯后来写道，在此期间的经历使他对“一切苏联官员、政客和委员们都产生了不可克服的生理性厌恶”。引自斯坦豪斯自述。

在两次大战之间以及苏联占领期间，斯坦豪斯向著名的《苏格兰笔记本》（Scottish Book）贡献了十个问题，其中包括最后一个——记录于1941年纳粹发动“巴巴罗萨行动”（Operation Barbarossa）攻占利沃夫前夕。\cite{4}

由于犹太血统，斯坦豪斯在纳粹占领期间被迫躲藏。他最初藏身于利沃夫的朋友家中，随后辗转至扎莫希奇（Zamość）附近的小镇奥西奇纳（Osiczyna），以及克拉科夫（Kraków）附近的贝尔德霍夫（Berdechów）。\cite{7,8}  
波兰反纳粹抵抗组织为他提供了一份伪造的身份证明，身份是早已去世的森林护林员“格热戈日·克罗赫马尔内”（\emph{Grzegorz Krochmalny}）。以此身份，他秘密教授课程（德军占领期间禁止波兰人接受高等教育）。  
由于担心一旦被捕便会被处死，斯坦豪斯在无书可查的情况下，仅凭记忆重建并记录下了自己掌握的全部数学知识，并撰写了大量回忆录——其中只有少部分在战后被发表。\cite{8}

此外，在藏匿期间，由于与外界隔绝、无法获得可靠的战争信息，斯坦豪斯设计了一种统计方法，通过分析地方报纸上零星刊登的讣告，来估算德军在前线的伤亡人数。  
该方法基于讣告中“死者为某人之子”、“某人之二子”、“某人之三子”等表述的相对频率。\cite{8}

据他的学生兼传记作者马克·卡茨（Mark Kac）回忆，斯坦豪斯曾说，自己一生中最幸福的一天，是“德军撤离波兰而苏军尚未来临的那二十四小时”——他说：“他们已经走了，而他们还没来。”\cite{8}
